\documentclass[10pt,a4paper]{amsart}
\usepackage[margin=19mm]{geometry}
\usepackage{amsmath,amssymb,mathtools}
\usepackage[scr=rsfs,cal=euler]{mathalfa}
\usepackage{xfrac}
\usepackage[unicode,hidelinks,bookmarks=false]{hyperref}
\newcommand{\ie}{i.e.\ }
\newcommand{\cf}{cf.\ }
\newcommand{\resp}{respectively\ }
\newcommand{\Subsection}{\subsection}
\providecommand{\nobreakdash}{\nobreak\hbox{-}}
\setlength{\emergencystretch}{2em}
\multlinegap0pt
\usepackage{bm}
\usepackage{bbm}
\usepackage{dsfont}
\usepackage{xspace}
\usepackage{cancel}
\usepackage{mathtools}
\usepackage{amsthm}
\makeatletter
\@ifundefined{thm}{\newtheorem{thm}{Theorem}}{}
\makeatother
\title[Waterborne epidemics via a new coupled SIR--Pathogen--Navier]{Waterborne epidemics via a new coupled SIR--Pathogen--Navier-Stokes system: Mathematical modeling, nonlinear analysis and numerical simulation}
\author{Mohamed Mehdaoui, Yassine Ouzrour}
\date{Source: arXiv:2511.04841v1 (CC BY 4.0)}
\begin{document}
\maketitle

\noindent\emph{Source and licence.} Waterborne epidemics via a new coupled SIR--Pathogen--Navier-Stokes system: Mathematical modeling, nonlinear analysis and numerical simulation. Version arXiv:2511.04841v1 (CC BY 4.0), distributed under the Creative Commons Attribution 4.0 International licence, as recorded in the arXiv licence metadata for this exact version and shown on the abstract page https://arxiv.org/abs/2511.04841v1. Full rights notice: licenses/arxiv-2511.04841-CC-BY-4.0.txt.\par\smallskip

\noindent\textit{Editorial scope.} Counted unit: the Thm whose source label is \texttt{None} in the article (source0.tex lines 887--889, source label \texttt{None}), delivered together with the article's own complete original proof of it (source0.tex lines 890--1094, one \texttt{proof} environment of 205 lines). No mathematics is written, repaired or re-derived: statement and proof are the article's own text line for line. Required context retained uncounted: SIRCNS (lines 249--262); Gagliardo-Nirenberg-U (lines 386--389); inq-Young (lines 392--394); eqvar1 (lines 446--463); eqvar2 (lines 466--475); galc (lines 729--731). Every definition, display and label of those lines is retained unchanged. Omitted from this package and not redistributed: 1--248, 263--385, 390--391, 395--445, 464--465, 476--728, 732--886, 1095--1432. Nothing inside a retained excerpt is omitted or altered: every retained body is delivered exactly as the article's lines are written, so no omission receipt of any kind accompanies this package. Citations: the retained excerpts carry no citation command at all, so no bibliography entry of the article is transcribed and no reference is redistributed. Reference adaptation: none was needed and none was made; every reference inside the delivered unit resolves inside it, so no unresolved reference or citation is produced. Printed slips kept literal: no printed word, symbol, hypothesis or number of the counted statement or of its proof was altered. Layout and class: the article's own class is \texttt{elsarticle} with packages including physics, amssymb, atbegshi, inputenc, alltt, lipsum, epstopdf, amsfonts, mathtools, amsmath, times, tikz-cd, manyfoot, booktabs; the wrapper is the lane's pinned \texttt{amsart} editorial wrapper at 10pt on A4, which restates the theorem-like environments the retained text needs and adds the article's own macro definitions that the retained text needs (none), with extra wrapper packages (bm, bbm, dsfont, xspace, cancel, mathtools). The delivered numbering is the unit's own. Assets: no retained span contains a figure environment, a graphics-inclusion command, a PDF-page inclusion, a table or a TikZ picture; no image file is copied into this package, the package declares no asset dependency and no image file is redistributed with it. Licence: version arXiv:2511.04841v1 of this article is distributed under the Creative Commons Attribution 4.0 International licence, as recorded in the arXiv licence metadata for this exact version and shown on the abstract page https://arxiv.org/abs/2511.04841v1. A full rights notice is delivered at licenses/arxiv-2511.04841-CC-BY-4.0.txt. Characters: every retained excerpt is pure ASCII, so the delivered engine, XeLaTeX with its default Latin Modern text font, reports no missing character.
	\begin{equation}\label{SIRCNS}
		\begin{cases}
			\begin{split}
				&\partial_t \boldsymbol{U} +(\boldsymbol{U}\cdot\nabla)\boldsymbol{U} = -\nabla p + \operatorname{div}(\nu(C)\nabla\boldsymbol{U} ) + \boldsymbol{f}, &&\text{in } Q_T,\\
				&\operatorname{div}(\boldsymbol{U})=0, &&\text{in } Q_T,\\
				&\partial_t C + \boldsymbol{U}\cdot\nabla C= D_C\Delta C + \alpha I  - \lambda C, &&\text{in } Q_T,\\
				&\partial_t S= D_S\Delta S - \beta(C)\frac{S I}{N}+\Lambda - \eta S, &&\text{in } Q_T,\\
				&\partial_t I= D_I\Delta I + \beta(C)\frac{S I}{N} - \gamma I - \eta I, &&\text{in } Q_T,\\
				&\partial_t R= D_R\Delta R + \gamma I- \eta R, &&\text{in } Q_T,\\
				&\boldsymbol{U}=C=\nabla S \cdot \overrightarrow{n}=\nabla I \cdot \overrightarrow{n}=\nabla R \cdot \overrightarrow{n}=0,  && \text{on } \Sigma_T,\\
				&(\boldsymbol{U}(.,0),C(.,0),S(.,0),I(.,0),R(.,0))=(\boldsymbol{U}_0,C_0,S_0,I_0,R_0) ,  && \text{in }\Omega.
			\end{split}
		\end{cases}
	\end{equation}

	\begin{align}
		\|\boldsymbol{Z}(t)\|_{\mathbf{L}^{4}(\Omega_b)} &\leq c\|\boldsymbol{Z}(t)\|_{\mathbf{H}^{1}(\Omega_b)}^{\zeta}\|\boldsymbol{Z}(t)\|_{\mathbf{L}^{2}(\Omega_b)}^{1-\zeta}, && \text{for all } \boldsymbol{Z} \in \mathbf{H}^{1}(\Omega_b), \label{Gagliardo-Nirenberg-U} \\
		\|\psi(t)\|_{L^{4}(\Omega)} &\leq c\|\psi(t)\|_{H^{1}(\Omega)}^{\zeta}\|\psi(t)\|_{L^{2}(\Omega)}^{1-\zeta}, && \text{for all } \psi \in H^{1}(\Omega) \label{Gagliardo-Nirenberg-C},
	\end{align}

	\begin{equation}\label{inq-Young}
		ab \leq \delta a^{p} + c(\delta) b^{q},
	\end{equation}

		\begin{eqnarray}\label{eqvar1}
			\begin{split}
				\int_{0}^{T}		\left\langle\partial_t\boldsymbol{U}, \boldsymbol{Z}\right\rangle\, dt+\iint_{Q_T}	\nu(C) \nabla \boldsymbol{U} : \nabla \boldsymbol{Z}\, d\mathbf{x}\, dt+\iint_{Q_T}	(\boldsymbol{U}\cdot\nabla)\boldsymbol{U} \cdot  \boldsymbol{Z}\, d\mathbf{x}\, dt\\
				-\iint_{Q_T}\boldsymbol{f}\cdot \boldsymbol{Z}\, d\mathbf{x}\, dt&=0,
				\\  
				\int_{0}^{T}	\left\langle\partial_tC, \psi_C \right\rangle\, dt+ \iint_{Q_T}  \boldsymbol{U}\cdot\nabla C\, \psi_C\, d\mathbf{x}\, dt+ D_C\iint_{Q_T}\nabla C \cdot\nabla \psi_C\, d\mathbf{x}\, dt\\ - \alpha\iint_{Q_T} I\, \psi_C\, d\mathbf{x}\, dt + \lambda \iint_{Q_T} C \,\psi_C \, d\mathbf{x}\, dt
				&=0,
				\\ 	  		
				\int_{0}^{T}	\left\langle\partial_t S, \psi_S \right\rangle\, dt+ D_S\iint_{Q_T}\nabla S \cdot\nabla \psi_S\, d\mathbf{x}\, dt+\iint_{Q_T}\beta(C) \frac{SI}{N} \, \psi_S\, d\mathbf{x}\, dt\\-\iint_{Q_T} \Lambda \,\psi_S \, d\mathbf{x}\, dt + \eta \iint_{Q_T} S \,\psi_S \, d\mathbf{x}\, dt
				&=0,
				\\ 	  		
				\int_{0}^{T}	\left\langle\partial_t I, \psi_I \right\rangle\,  dt+ D_I\iint_{Q_T}\nabla I \cdot\nabla \psi_I\, d\mathbf{x}\, dt-\iint_{Q_T}\beta(C) \frac{SI}{N} \, \psi_I\, d\mathbf{x}\, dt + (\eta+\gamma) \iint_{Q_T} I \,\psi_I \, d\mathbf{x}\, dt
				&=0,
				\\ 	  		
			\int_{0}^{T}	\left\langle\partial_t R, \psi_R \right\rangle\,  dt+ D_R\iint_{Q_T}\nabla R \cdot\nabla \psi_R\, d\mathbf{x}\, dt-\gamma\iint_{Q_T} I \,\psi_R \, d\mathbf{x}\, dt+ \eta\iint_{Q_T} R \,\psi_R \, d\mathbf{x}\, dt
				&=0,
			\end{split} 
		\end{eqnarray}   

		\begin{eqnarray}\label{eqvar2}
			\begin{split}
				\boldsymbol{U}(0,\mathbf{.}) &=\boldsymbol{U}_{0},  & \hbox { in }   \Omega, \\
				C(0,\mathbf{.}) &=C_{0},  & \hbox { in }   \Omega,
				\\
				S(0,\mathbf{.}) &=S_{0},  & \hbox { in }   \Omega, \\
				I(0,\mathbf{.}) &=I_{0},  & \hbox { in }   \Omega, \\
				R(0,\mathbf{.}) &=R_{0},  & \hbox { in }   \Omega,  
			\end{split}
		\end{eqnarray}

		\begin{equation}\label{galc}
\displaystyle \int_\Omega  \boldsymbol{U}^n\cdot\nabla C^n \, C^n \, d\mathbf{x}= \frac{1}{2}\int_\Omega\boldsymbol{U}^n\nabla ((C^n)^2) \, d\mathbf{x}=-\frac{1}{2}\int_\Omega\operatorname{div}(\boldsymbol{U}^n) \, (C^n)^2 \, d\mathbf{x}+\frac{1}{2}\int_{\partial\Omega} \boldsymbol{U}^n\cdot\boldsymbol{n}\,(C^n)^2\,ds=0.
		\end{equation}%\mathbf{U}

\begin{thm}
		The weak solution of (\ref{SIRCNS}) is unique.
\end{thm}

	\begin{proof}
		Suppose that there are two solutions $$(\boldsymbol{U}_1,  C_1, S_1, I_1, R_1)$$ and 
		$$(\boldsymbol{U}_2, C_2, S_2, I_2, R_2)$$ 
		to the problem (\ref{eqvar1})-(\ref{eqvar2}). 
		
		\noindent
{Denote $\boldsymbol{U}:= \boldsymbol{U}_1 - \boldsymbol{U}_2$,  $C:= C_1 - C_2$, $S:= S_1 - S_2$, $I:= I_1 - I_2$ and $R:= R_1 - R_2$.   
		 Then, $\boldsymbol{U}$, $C$, $S$, $I$ and  $R$ satisfy the following equations}
		\begin{eqnarray}
			\left\langle\partial_{t}\boldsymbol{U}, \boldsymbol{Z}\right\rangle+ \int_{\Omega} \nu(C_1) \nabla \boldsymbol{U}: \nabla \boldsymbol{Z}~d{\mathbf{x}}+\int_{\Omega}\left[\nu(C_1)-\nu\left(C_{2}\right)\right] \mathbb{D}(\boldsymbol{U}_{2}): \mathbb{D}(\boldsymbol{Z}) ~d{\mathbf{x}}&
			\nonumber\\
			+\int_{\Omega}( \boldsymbol{U}\cdot \nabla) \boldsymbol{U}_{2}\cdot \boldsymbol{Z}~d{\mathbf{x}}+\int_{\Omega}( \boldsymbol{U}_1\cdot \nabla) \boldsymbol{U}\cdot \boldsymbol{Z}~d{\mathbf{x}}
			&=0, \;\;\label{uni1} \\  
			\left\langle\partial_tC, \psi_C\right\rangle+D_C\int_{\Omega}\nabla C\cdot\nabla\, \psi_C~d{\mathbf{x}}+\int_{\Omega}\boldsymbol{U}\cdot\nabla C_1 \psi_C~d{\mathbf{x}}+\int_{\Omega} \boldsymbol{U}_{2} \cdot\nabla C\, \psi_C~d{\mathbf{x}}
			&
			\nonumber\\
			-\alpha\int_{\Omega} I\, \psi_C~d{\mathbf{x}}+\lambda\int_{\Omega} C\, \psi_C~d{\mathbf{x}}
			&=0,\label{uni2} \\ 
			\left\langle\partial_tS, \psi_S\right\rangle+D_S\int_{\Omega}\nabla S\cdot\nabla\, \psi_S ~d{\mathbf{x}}+\int_{\Omega}\Big[\beta(C_1)\frac{S_1 I_1}{N_1}-\beta(C_2)\frac{S_2 I_2}{N_2}\Big]\psi_S~d{\mathbf{x}}
			%		&
			%		\nonumber\\
			+\eta\int_{\Omega} S\, \psi_S~d{\mathbf{x}}&= 0, \;\;\label{uni3}
			\\ 
			\left\langle\partial_t I, \psi_I\right\rangle+D_I\int_{\Omega}\nabla I\cdot\nabla\, \psi_I ~d{\mathbf{x}}-\int_{\Omega}\Big[\beta(C_1)\frac{S_1 I_1}{N_1}-\beta(C_2)\frac{S_2 I_2}{N_2}\Big]\psi_I~d{\mathbf{x}}
			%		&
			%		\nonumber\\
			+(\eta+\gamma)\int_{\Omega} I\, \psi_I~d{\mathbf{x}}&= 0, \;\;\label{uni4}
			\\ 
			\left\langle\partial_tR, \psi_R\right\rangle+D_R\int_{\Omega}\nabla R\cdot\nabla\, \psi_R ~d{\mathbf{x}}
			+\eta\int_{\Omega} R\, \psi_R~d{\mathbf{x}}-\gamma\int_{\Omega} I\, \psi_R~d{\mathbf{x}}&= 0, \;\;\label{uni5}
		\end{eqnarray} 
		for all test functions  
		$\boldsymbol{Z} \in D([0,T); \mathbf{H})$, $\psi_C \in D([0,T);  H^1_0)$, $\psi_S \in D([0,T);  H^1_0)$, $\psi_I \in D([0,T);  H^1_0)$ 
		and $\psi_R \in D([0,T);  H^1_0)$,  
		with the initial condition
		$$\begin{aligned}
			\boldsymbol{U}(0,\mathbf{.})=C(0,\mathbf{.})=S(0,\mathbf{.})=I(0,\mathbf{.})=R(0,\mathbf{.})&=0& \text{ in } & \Omega.
			\end{aligned}
		$$
		Now, by using $\boldsymbol{Z} = \boldsymbol{U}(t)$ as a test function in $(\ref{uni1})$, we acquire the following estimate
		\begin{equation}\label{testU}
			\begin{aligned}
				\frac{1}{2} \frac{\mathrm{d}}{\mathrm{d} t}\|\boldsymbol{U}(t)\|_{\boldsymbol{L}^2(\Omega)}^{2}+\nu_{1}\|\boldsymbol{U}(t)\|_{\boldsymbol{H}^1(\Omega)}^{2} & \leq \left|\int_{\Omega}( \boldsymbol{U}(t)\cdot \nabla) \boldsymbol{U}_{2}(t)\cdot \boldsymbol{U}(t)~d{\mathbf{x}}\right|+\left|\int_{\Omega}( \boldsymbol{U}_1(t)\cdot \nabla) \boldsymbol{U}(t)\cdot \boldsymbol{U}(t)~d{\mathbf{x}}\right|
				\\
				& \hspace{1  cm}
				+\left|\int_{\Omega}(\nu(C_1)-\nu(C_{2})) \mathbb{D}(\boldsymbol{U}_{2}(t)): \mathbb{D}(\boldsymbol{U}(t))~d{\mathbf{x}}\right|.
			\end{aligned} 
		\end{equation}
		Based on (\ref{Gagliardo-Nirenberg-U}) and (\ref{inq-Young}), the first term in (\ref{testU}) can be estimate as follows
		\begin{equation}\label{estem-b}
			\begin{split} 
				\hspace*{-0.2cm}
				\left|\int_{\Omega}( \boldsymbol{U}(t)\cdot \nabla) \boldsymbol{U}_{2}(t)\cdot \boldsymbol{U}(t)~d{\mathbf{x}}\right| 
				&\leq {\|\boldsymbol{U}(t)\|}_{\boldsymbol{L}^4(\Omega)} {\|\nabla\boldsymbol{U}_2(t)\|}_{\boldsymbol{L}^2(\Omega)} {\|\boldsymbol{U}(t)\|}_{\boldsymbol{L}^4(\Omega)}
				\\
				&\leq c{\|\boldsymbol{U}(t)\|}_{\boldsymbol{H}^1(\Omega)}^{1/2}
				{\|\boldsymbol{U}(t)\|}_{\boldsymbol{L}^2(\Omega)}^{1/2} \left\|\boldsymbol{U}_{2}(t)\right\|_{\mathbf{H}^1(\Omega)}{\|\boldsymbol{U}(t)\|}_{\boldsymbol{H}^1(\Omega)}^{1/2}
				{\|\boldsymbol{U}(t)\|}_{\boldsymbol{L}^2(\Omega)}^{1/2}
				\\
				&\leq \delta\|\boldsymbol{U}(t)\|_{\boldsymbol{H}^1(\Omega)}^{2}+c_{\delta}\left\|\boldsymbol{U}_{2}(t)\right\|_{\mathbf{H}^{1}(\Omega)}^2\|\boldsymbol{U}(t)\|_{\mathbf{L}^{2}(\Omega)}^{2}.
			\end{split}
		\end{equation}
		For the second term, by Green's formula, we have that
		\begin{equation}\label{estm-b-null} 
			\int_{\Omega}( \boldsymbol{U}_1(t)\cdot \nabla) \boldsymbol{U}(t)\cdot \boldsymbol{U}(t)~d{\mathbf{x}}=0.
		\end{equation}
		In addition, the last term in (\ref{testU})  can be estimates by means of  Young inequality as well as H{\"o}lder's inequality to obtain
		\begin{equation}\label{estmD}
			\begin{aligned}
				\left|\int_{\Omega}\left(\nu(C_1)-\nu(C_{2})\right) \mathbb{D}(\boldsymbol{U}_{2}): \mathbb{D}(\boldsymbol{U})~d{\mathbf{x}}\right|
				& \leq \|(\nu(C_1)-\nu(C_{2}))\|_{L^{\infty}(\Omega)}
				\| \mathbb{D}(\boldsymbol{U}_{2})\|_{\mathbf{L}^{2}(\Omega)} \|\mathbb{D}(\boldsymbol{U})\|_{\mathbf{L}^{2}(\Omega)} \\
				& \leq c\|C(t)\|_{L^2(\Omega)}
				\| \mathbb{D}(\boldsymbol{U}_{2})\|_{\boldsymbol{L}^{2}(\Omega)} \|\boldsymbol{U}(t)\|_{\mathbf{H}^{1}(\Omega)}\\
				& \leq \delta \|\boldsymbol{U}(t)\|_{\boldsymbol{H}^{1}(\Omega)}^2+C_\delta
				\|\boldsymbol{U}_{2}(t)\|_{\mathbf{H}^{1}(\Omega)}^2\|C(t)\|_{L^2(\Omega)}^2.
			\end{aligned}
		\end{equation}
		Consequently, the estimates $(\ref{estem-b})-(\ref{estmD})$ imply that
		\begin{equation}\label{estmU}
			\begin{aligned}
				\frac{\mathrm{d}}{\mathrm{d} t}\|\boldsymbol{U}(t)\|_{\boldsymbol{L}^2(\Omega)}^{2}+c\|\boldsymbol{U}(t)\|_{\boldsymbol{H}^1(\Omega)}^{2}  & \leq \delta \|\boldsymbol{U}(t)\|_{\boldsymbol{H}^{1}(\Omega)}^2+
				c_{\delta}\|\boldsymbol{U}_{2}(t)\|_{\mathbf{H}^{1}(\Omega)}^2\left(\|\boldsymbol{U}(t)\|_{\mathbf{L}^{2}(\Omega)}^{2}+\|C(t)\|_{L^2(\Omega)}^2\right), 
			\end{aligned} 
		\end{equation} 
		%		where $R_1(t)=\Big(\left\|\boldsymbol{U}_{1}(t)\right\|_{\mathbf{L}^{4}}^{\frac{2}{1-\zeta}}+\left\|\boldsymbol{U}_{2}(t)\right\|_{\mathbf{w}^{1,2}}^{\frac{1}{1-\zeta}}+\| \mathbb{D}(\boldsymbol{U}_{2})\|_{\mathbf{L}^{4}}^2+1\Big)$.\\ %and $\beta_1(t)=C(\delta)\left(\left\|\boldsymbol{U}_{1}(t)\right\|_{\mathbf{L}^{4}}^{\frac{2}{1-\zeta}}+\left\|\boldsymbol{U}_{2}(t)\right\|_{\mathbf{w}^{1,2}}^{\frac{1}{1-\zeta}}+1\right)$.\\
		Now, we substitute $\psi_C=C$ in (\ref{uni2}), to obtain the
		following estimate   
		$$
		\begin{array}{rl}
			\displaystyle \frac{1}{2}\frac{\mathrm{d}}{\mathrm{d} t}\|C(t)\|_{L^2(\Omega)}^{2}+D_C\left\|C\right\|_{H^{1}(\Omega)}^{2}&\leq \displaystyle \Big|\int_{\Omega}\boldsymbol{U}\cdot\nabla C_1 \,C~d{\mathbf{x}}\Big| +\Big|\int_{\Omega} \boldsymbol{U}_{2} \cdot\nabla C\, C~d{\mathbf{x}}\Big| \\&\displaystyle\qquad+\alpha\left| \int_{\Omega} I\, C~d{\mathbf{x}} \right|+\lambda\left|\int_{\Omega} C^2~d{\mathbf{x}} \right|.
		\end{array}
		$$
		By virtue of  H{\"o}lder and Young inequalities, we estimate the terms on the right-hand side of the previous equation as follows
		
		\begin{equation}\label{C1-C2}
			\Big|\int_{\Omega} I C ~d{\mathbf{x}} \Big| \leq \frac{1}{2}\Big(\left\|I\right\|_{L^{2}(\Omega)}^{2}+\left\|C\right\|_{L^{2}(\Omega)}^{2}\Big).
		\end{equation}
		By the Gagliardo-Nirenberg inequality with $d=2$, we deduce that
		\begin{equation}
			\begin{aligned}
				&\Big|\int_{\Omega}\boldsymbol{U}(t)\cdot\nabla C_1(t) \,C(t)~d{\mathbf{x}}\Big|
				\\ &\qquad \leq
				\|\boldsymbol{U}(t)\|_{\mathbf{L}^{4}(\Omega)}\left\|\nabla C_1(t)\right\|_{\mathbf{L}^{2}(\Omega)}\left\|C(t)\right\|_{L^{4}(\Omega)} 
				\\ &\qquad  \leq
				c\|\boldsymbol{U}(t)\|_{\mathbf{H}^{1}(\Omega)}^{1/2}\|\boldsymbol{U}(t)\|_{\mathbf{L}^{2}(\Omega)}^{1/2}\left\|C_1(t)\right\|_{H^{1}(\Omega)}^{}\left\|C(t)\right\|_{H^{1}(\Omega)}^{1/2}\left\|C\right\|_{L^{2}(\Omega)}^{1/2}
				\\ &\qquad  \leq
				\delta(\left\|C(t)\right\|_{H^{1}(\Omega)}^{}\|\boldsymbol{U}(t)\|_{\mathbf{H}^{1}(\Omega)})+c_{\delta}\left\|C_1\right\|_{H^{1}(\Omega)}^{2}\|\boldsymbol{U}(t)\|_{\mathbf{L}^{2}(\Omega)}^{}\left\|C(t)\right\|_{L^2(\Omega)}
				\\ &\qquad \leq
				\delta/2\Big(\left\|C(t)\right\|_{H^{1}(\Omega)}^{2}\hspace*{-0.1 cm}+\|\boldsymbol{U}(t)\|_{\mathbf{H}^{1}(\Omega)}^{2}\Big)\hspace*{-0.1 cm}+c_{\delta}\left\|C_1\right\|_{H^{1}(\Omega)}^{2}\hspace*{-0.1 cm}\Big(\|\boldsymbol{U}(t)\|_{\mathbf{L}^{2}(\Omega)}^{2}\hspace*{-0.1 cm}+\left\|C(t)\right\|_{L^2(\Omega)}^2\Big),
			\end{aligned}
		\end{equation}
	for a given positive function $c_{\delta}$ depending on $\delta$.
	
	
	\noindent
		Moreover, using the same technique in (\ref{galc}), we obtain
		\begin{equation}\label{E2}
			\begin{aligned}
				\int_{\Omega} \boldsymbol{U}_{2} \cdot\nabla C\, C~d{\mathbf{x}}=0.
			\end{aligned}
		\end{equation} 
		Collecting the previous results \eqref{C1-C2}-\eqref{E2}, we deduce that
		\begin{equation}\label{estmtheta}
			\begin{split}
				\frac{d}{d t}\left\|C\right\|_{L^{2}(\Omega)}^{2}&+c\left\|C(t)\right\|_{H^{1}(\Omega)}^{2}
				\\&\leq \delta\left(\left\|C(t)\right\|_{H^{1}(\Omega)}^{2}+\|\boldsymbol{U}\|_{\mathbf{H}^{1}(\Omega)}^{2}\right)+c_{\delta}\left\|C_1\right\|_{H^{1}(\Omega)}^{2} \left(\|C\|_{L^{2}(\Omega)}^{2}+\|\boldsymbol{U}\|_{\mathbf{L}^{2}(\Omega)}^{2}\right)\\
				&\qquad+\frac{1}{2}\Big(\left\|I\right\|_{L^{2}(\Omega)}^{2}+\left\|C\right\|_{L^{2}(\Omega)}^{2}\Big)+\lambda\left\|C\right\|_{L^{2}(\Omega)}^{2}.
			\end{split}
		\end{equation}
		
		\noindent For $\psi_S=S$ in (\ref{uni3}), we obtain that 
		\begin{equation}\label{estmS0}
			\displaystyle \frac{1}{2}\frac{\mathrm{d}}{\mathrm{d} t}\|S(t)\|_{L^2(\Omega)}^{2}+D_C\left\|S(t)\right\|_{H^{1}(\Omega)}^{2}\leq  \Big|\int_{\Omega}\Big[\beta(C_1)\frac{S_1 I_1}{N_1}-\beta(C_2)\frac{S_2 I_2}{N_2}\Big] S~d{\mathbf{x}}\Big| +\eta\left|\int_{\Omega} S^2~d{\mathbf{x}} \right|.
		\end{equation}
		We get that
		%{\color{red}	
			%Si la fonction	$(x,y)\longrightarrow \frac{x y}{x+y}$,	for $x,y>0$ est elle satisfais que $$ \Big|\frac{x_1 y_1}{x_1+y_1}-\frac{x_2 x_2}{x_2+y_2}\Big|\leq L_c(|x_1-x_2|+|y_1-y_2|) $$
			%Dans ce cas on a:
			%}
		\begin{equation*}
			\begin{split}
				\Big|\int_{\Omega}\Big[\beta(C_1)\frac{S_1 I_1}{N_1}&-\beta(C_2)\frac{S_2 I_2}{N_2}\Big] S~d{\mathbf{x}}\Big|
				\\&=\Big|\int_{\Omega}\beta(C_1)\Big[\frac{S_1 I_1}{N_1}-\frac{S_2 I_2}{N_2}\Big] S~d{\mathbf{x}}+\int_{\Omega}\Big[\beta(C_1)-\beta(C_2)\Big]\frac{S_2 I_2}{N_2} S~d{\mathbf{x}}\Big|
				\\&\leq \beta_2\int_{\Omega}\Big|\frac{S_1 I_1}{N_1}-\frac{S_2 I_2}{N_2}\Big| S~d{\mathbf{x}}+\int_{\Omega}\Big|\beta(C_1)-\beta(C_2)\Big|\frac{S_2 I_2}{N_2} S~d{\mathbf{x}}\Big|
				\\&\leq\beta_2\left\|\frac{S_1 I_1}{N_1}-\frac{S_2 I_2}{N_2}\right\|_{L^{2}(\Omega)}\left\|S\right\|_{L^{2}(\Omega)}+ \left\|\beta(C_1)-\beta(C_2)\right\|_{L^{4}(\Omega)} \left\|\frac{S_2 I_2}{N_2}\right\|_{L^{2}(\Omega)} \|S\|_{L^{4}(\Omega)}
				\\&\leq \beta_2 L_g(\left\|S\right\|_{L^{2}(\Omega)}+\left\|I\right\|_{L^{2}(\Omega)})\|S\|_{L^{2}(\Omega)}+ L_{\beta}\|C\|_{L^{4}(\Omega)}\left\|S_2\right\|_{L^{2}(\Omega)} \|S\|_{L^{4}(\Omega)}
				\\&\leq \beta_2 L_g(\left\|S\right\|_{L^{2}(\Omega)}+\left\|I\right\|_{L^{2}(\Omega)})\|C\|_{L^{2}(\Omega)}
				\\
			    &+ c\left\|C(t)\right\|_{H^{1}(\Omega)}^{1/2}\left\|C(t)\right\|_{L^{2}(\Omega)}^{1/2}\left\|S_2\right\|_{L^{2}(\Omega)} \left\|S(t)\right\|_{H^{1}(\Omega)}^{1/2}\left\|S(t)\right\|_{L^{2}(\Omega)}^{1/2}
				\\&\leq \delta \left(\left\|S\right\|_{H^{1}(\Omega)}^2+\left\|C\right\|_{H^{1}(\Omega)}^2\right)+c\left( \left\|C\right\|_{L^{2}(\Omega)}^2+\left\|I\right\|_{L^{2}(\Omega)}^2+\left\|S\right\|_{L^{2}(\Omega)}^2\right)
				\\&+c_{\delta}\left\|S_2\right\|_{H^{1}(\Omega)}^{2}\left( \left\|C\right\|_{L^{2}(\Omega)}^2+\left\|S\right\|_{L^{2}(\Omega)}^2\right). 
			\end{split}
		\end{equation*}
		Thereby
		\begin{equation}\label{estmS}
			\begin{split}
				&\frac{\mathrm{d}}{\mathrm{d} t}\|S(t)\|_{L^2(\Omega)}^{2}+c_1\left\|S\right\|_{H^{1}(\Omega)}^{2}\\
				&\quad\leq \delta \left(\left\|S\right\|_{H^{1}(\Omega)}^2+\left\|C\right\|_{H^{1}(\Omega)}^2\right)+\left(c_{\delta}\left\|S_2\right\|_{H^{1}(\Omega)}^{\frac{1}{1-\zeta}}+c_2\right)\left( \left\|C\right\|_{L^{2}(\Omega)}^2+\left\|I\right\|_{L^{2}(\Omega)}^2+\left\|S\right\|_{L^{2}(\Omega)}^2\right). 
			\end{split}
		\end{equation}
		
		
		\noindent Similarly to (\ref{estmS}),  we use $\psi_I=I$ as a test function in (\ref{uni4}) to deduce that   
		\begin{equation}\label{estmI}
			\begin{split}
				&\displaystyle \frac{\mathrm{d}}{\mathrm{d} t}\|I(t)\|_{L^2(\Omega)}^{2}+c_3\left\|I\right\|_{H^{1}(\Omega)}^{2}
				\\
				&\quad\leq  \Big|\int_{\Omega}\Big[\beta(C_1)\frac{S_1 I_1}{N_1}-\beta(C_2)\frac{S_2 I_2}{N_2}\Big] I~d{\mathbf{x}}\Big| +(\eta+\gamma)\left|\int_{\Omega} I^2~d{\mathbf{x}} \right|
				\\
				&\quad\leq \delta \left(\left\|I\right\|_{H^{1}(\Omega)}^2+\left\|C\right\|_{H^{1}(\Omega)}^2\right)+\left(c_{\delta}\left\|S_2\right\|_{H^{1}(\Omega)}^{\frac{1}{1-\zeta}}+c_4\right)
%				\\
%				&\times 
				\left( \left\|C\right\|_{L^{2}(\Omega)}^2+\left\|I\right\|_{L^{2}(\Omega)}^2+\left\|S\right\|_{L^{2}(\Omega)}^2\right) 
				.
			\end{split}
		\end{equation}
		
		
		\noindent Further, we use substitute $\psi_R=R$ in (\ref{uni5}) to get   
		\begin{equation}\label{estmR}
			\begin{split}
				\displaystyle \frac{1}{2}\frac{\mathrm{d}}{\mathrm{d} t}\|R(t)\|_{L^2(\Omega)}^{2}+D_R\left\|R\right\|_{H^{1}(\Omega)}^{2}&\leq  \gamma \Big|\int_{\Omega}I\, R~d{\mathbf{x}}\Big| +\eta\left|\int_{\Omega} R^2~d{\mathbf{x}} \right|
				\\&\leq\gamma\left\|I\right\|_{L^{2}(\Omega)}\left\|C\right\|_{L^{2}(\Omega)}+\eta \left\|C\right\|_{L^{2}(\Omega)}^{2}
				\\&\leq\max(\gamma/2,\eta+\gamma/2)\Big(\left\|I\right\|_{L^{2}(\Omega)}^2+ \left\|C\right\|_{L^{2}(\Omega)}^{2}\Big)
				.
			\end{split}
		\end{equation}
		
		
		\noindent	We sum up the estimates (\ref{estmU}), (\ref{estmtheta}),  (\ref{estmS}), (\ref{estmI}) and (\ref{estmR}) to deduce, for a sufficiently small $\delta>0$
		\begin{equation}\label{estm-U-Theta}
			\begin{aligned}
				\frac{d}{d t}\Big(\|\boldsymbol{U}\|_{\mathbf{L}^{2}}^{2}&+ \left\|C\right\|_{L^{2}(\Omega)}^{2}+\left\|S\right\|_{L^{2}(\Omega)}^{2}+\left\|I\right\|_{L^{2}(\Omega)}^{2}+\left\|R\right\|_{L^{2}(\Omega)}^{2} \Big)\\ & \leq K(t)\left(\|\boldsymbol{U}\|_{\mathbf{L}^{2}}^{2}+ \left\|C\right\|_{L^{2}(\Omega)}^{2}+\left\|S\right\|_{L^{2}(\Omega)}^{2}+\left\|I\right\|_{L^{2}(\Omega)}^{2}+\left\|R\right\|_{L^{2}(\Omega)}^{2}  \right),
			\end{aligned} 
		\end{equation}
		where $$K(t)=c_{\delta}\left(\left\|S_2(t)\right\|_{H^{1}(\Omega)}^{2}+\left\|C_1(t)\right\|_{H^{1}(\Omega)}^{2}+\|\boldsymbol{U}_{2}(t)\|_{\mathbf{H}^{1}(\Omega)}^2+1\right).$$
		
		\noindent
		Applying Gronwall's inequality to $(\ref{estm-U-Theta})$ along with the fact that $$\boldsymbol{U}(0,\mathbf{.})=C(0,\mathbf{.})=S(0,\mathbf{.})=I(0,\mathbf{.})=R(0,\mathbf{.})=0$$  
		we arrive at
		$$C=S=I=R=\boldsymbol{U}=0.$$
		\noindent
		This yields the uniqueness of the weak solution. 
	\end{proof}

\end{document}
